\scanpage{126}

Dans le cas exceptionnel de la double pyramide, on a $\Pi$ réduit à un élément \unsure{$A_0$|$F_0$|$P_0$} \deleted{\unsure{et}} i.e. $Z'=\{1,-\}$ sss $n$ impair ou $n=4$, et dans le cas inverse ($n$ pair $\geqslant6$) $\Pi$ a exactement deux éléments, qui correspondent à la double pyramide sur le = ensemble $E$ ; muni des deux structures polygonales différentes (d'ordre \deleted{pair} \inserted{$\deleted{\unsure{n}}n\equiv0\ (4)$ et} $\geqslant8$) déduites l'une de l'autre par ``multiplication par $\frac{n}{2}-1$''\ldots

\noindent\begin{tabular}{|p{.92\linewidth}}
\unsure{$\vert$|!}\quad Je laisse tomber par la suite ce cas de la double pyramide, sauf dans le cas de l'octaèdre, i.e pratiquement je demande que les \deleted{points}\inserted{sommets} aient tous même ordre -- donc le seul cas exceptionnel où $\Pi$ n'est pas un torseur sous \deleted{\unsure{$Z$}} $Z'$ est alors celui où $P_0$ est l'octaèdre, $\Pi$ réduit à un élément, $Z'=\{1,\underline-\}$. 
\end{tabular}

Notons que par transport de structure, toute permutation $f$ de $S$ qui transforme \uline{un} élément $F_0$ de $\Pi$ en un autre $F_1\in\Pi$, transforme \inserted{\unsure{=}} $\Pi$ \inserted{tout entier} en lui-\deleted{\unsure{=}}même, et induit une bijection de $\Pi$ avec lui =. Soit donc
% archive sheet number: 125
